\section{Introduction}
\label{sec:intro}

\IEEEPARstart{C}{ontent} delivery networks (CDNs) protect paywalled video, software downloads and other private objects with \emph{access tokens}. An issuer, usually the content owner's login service, gives an authorized client a short-lived token. The token binds a URL prefix, an expiry time and optional client attributes. The client attaches it to every request: as a query parameter (a \emph{signed URL}), as a cookie, or in a header. Any edge server can check the token before serving from cache, without contacting the origin. How many requests one token carries varies widely. A token scoped to a streaming session covers one manifest request and then hundreds of segment requests. A signed URL for a single object may be used only once or twice.

Deployed schemes use one of two approaches. Symmetric schemes such as Akamai's EdgeAuth~\cite{akamai-edgeauth} append an HMAC~\cite{rfc2104} under a key held by every edge. Public-key schemes sign the token: CloudFront key groups accept RSA-2048 or ECDSA P-256 keys~\cite{cloudfront-signers}, and Google Media CDN uses Ed25519~\cite{mediacdn-sig}. HMAC-SHA256 is not threatened by quantum computers, but every verifier can also mint tokens. CDNs already hold many customers' secrets~\cite{liang2014httpscdn,cangialosi2016keysharing}, so a single compromised edge exposes every path the key covers. The public-key schemes keep edges secret-free, but Shor's algorithm~\cite{shor1997} would recover their signing keys from long-lived public keys. Authentication has no harvest-now-decrypt-later exposure. The reason to migrate now is lead time: token formats, SDKs and edge verifiers must change before a cryptographically relevant quantum computer (CRQC) exists, and industry roadmaps target post-quantum authentication by about 2029~\cite{cf2026roadmap}.

The obvious replacement is one FIPS~204 ML-DSA signature per token~\cite{fips204,rfc9964}. It produces a \num{3350}-B token, which every request then repeats. Two remedies are available. \emph{Stapling} detaches the signature and sends it once per connection (M2). \emph{Merkle batch signing}~\cite{merkle1989,batchsig2024,fregly2023mtl} signs one root for many tokens, which our earlier version of this work applied to CDN tokens as EMBER. Review of that version showed that at the operating point we had recommended (per-tenant trees, 10~ms epochs) most trees hold a single token and most of the edge benefit came from detaching the signature. A replay of Meta's public CDN traces~\cite{berg2020cachelib,cachelib-traces} then showed that even this edge verdict holds only when tokens carry many requests. This paper separates the two mechanisms, measures each against strong alternatives, and identifies what decides between them at the edge: the number of requests $K$ a token carries at one verifier.

\emph{Contributions.}
\begin{enumerate}
\item \textbf{A trust and threat comparison} of 13 baselines and M2 (Section~\ref{sec:threat}, Table~\ref{tab:trust}). It covers who can mint tokens after an edge compromise, the post-quantum assumption and token size; HMAC is already post-quantum secure.
\item \textbf{Pooled batching relieves the issuer} (Sections~\ref{sec:res-oprg} and~\ref{sec:res-trace}). EMBER, a CDN profile of MTL-mode and batch signatures rather than a new primitive, pays off when one tree is pooled per issuer key, which the pooled tenants must share, and the pool exceeds about $10^3$ tokens/s at $\Delta=10$~ms. It cuts issuer CPU 19--144$\times$ in simulation. At the full rates implied by the trace, per-content-type trees hold a median of 547 tokens and the issuer needs 0.38--1.08 cores against 177 for per-token signing; the single-token trees we reported earlier appear only at low rates.
\item \textbf{Requests per token decide the edge} (Sections~\ref{sec:res-edge} and~\ref{sec:res-trace}). With session tokens ($K\approx150$), M2 matches HMAC in whole-request edge CPU in a Go prototype and is the smallest post-quantum option on the wire. With one token per object per hour ($K=1.7$--2.4 on the trace), M2 verifies a signature about every second request and pooled EMBER is 28$\times$ cheaper in edge crypto and 4.4$\times$ smaller on the wire. From measured costs we derive the crossover: $K^*\approx48$--154 for CPU and 9--27 for bytes.
\item \textbf{HPACK behavior of real clients} (Section~\ref{sec:res-wire}). Chrome~153 adopts an edge-advertised 64~KiB HPACK table, so a 3.3~KB ML-DSA cookie costs 25.7~B per request; curl/nghttp2 and Go clamp it to 4096~B, yet Go kept the cookie indexed (228~B per request) while curl re-sent it.
\item \textbf{A prototype and a simulator.} The simulator matches the prototype on miss and byte counts in all 12 configurations, and we report where it fails (Section~\ref{sec:res-validate}). The prototype runs M2 with its retry path and pooled EMBER for edge CPU, latency under load, and recovery after cache loss, restarts and undersized caches (Sections~\ref{sec:design} and~\ref{sec:res-edge}).
\item \textbf{Deployment guidance} keyed to $K$, the per-key token rate and who holds the issuer key (Section~\ref{sec:discussion}), including when delegated Ed25519 keys suffice: only if a key cannot be broken within its exposure window, for which published estimates range from 38~h to 10~min per key.
\end{enumerate}
